\subsection{展开到基b}

命题8. 设 \(b\) 是一个整数 \(> 1\) 。对于每个整数 \(k > 0\) ，令 \(\mathbf{E}_k\) 为族 \((\mathrm{J}_h)_{0 \leqslant h \leqslant k-1}\) （§2，第6号）的字典序乘积，其中所有区间都与 {[}0, \(b - 1\) {]} 相同。对于每一个 \(r = (r_0, r_1, \ldots, r_{k-1}) \in \mathbf{E}_k\) ，令 \(f_k(r) = \sum_{h=0}^{k-1} r_h b^{k-h-1}\) ；则映射 \(f_k\) 是序集 \(\mathbf{E}_k\) 到区间 {[}0, \(b^k - 1\) {]} 上的一个严格递增同构.

证明通过对 \(k\) 进行归纳。对于 \(k = 1\) ，这是定义的直接推论。对于每个 \(r = (r_0, \ldots, r_{k-1}, r_k) \in \mathbf{E}_{k+1}\) ，令

\[
\varphi (r) = \left(r _ {0}, \dots , r _ {k - 1}\right) \in \mathrm{E} _ {k}.
\]

那么映射 \(r \to (\varphi(r), r_k)\) 是 \(\mathbf{E}_{k+1}\) 到 \(\mathbf{E}_k\) 与 \(J = [0, b - 1]\) 的字典序乘积上的一个同构；这由定义直接得出。我们可以写

\[
f _ {k + 1} (r) = b \cdot f _ {k} (\varphi (r)) + r _ {k};
\]

让我们证明关系 \(r < r'\) 在 \(\mathbf{E}_{k+1}\) 中蕴含 \(f_{k+1}(r) < f_{k+1}(r')\) 。事实上，我们或者有 \(\varphi(r) < \varphi(r')\) ，或者 \(\varphi(r) = \varphi(r')\) 并且 \(r_k < r'_k\) . 在第一种情形下，归纳假设蕴含 \(f_k(\varphi(r)) < f_k(\varphi(r'))\) ，因此（§4，第2号，命题2） \(f_k(q(r')) \geqslant f_k(\varphi(r)) + 1\) ；因此

\[
f _ {k - 1} (r ^ {\prime}) \geqslant b. f _ {k} (\varphi (r)) + b > f _ {k + 1} (r),
\]

因为 \(r_k \leqslant b - 1\) （第2号，命题3）。另一方面，如果 \(\varphi(r) = \varphi(r')\) 并且 \(r_k < r_k'\) ，显然 \(f_{k+1}(r) < f_{k+1}(r')\) 。现在归纳假设表明 \(f_k(\varphi(r)) \leqslant b^k - 1\) ，由此

\[
f _ {k + 1} (r) \leqslant b (b ^ {k} - 1) + b - 1 = b ^ {k + 1} - 1.
\]

由此得出 \(f_{k+1}\) 是 \(\mathbf{E}_{k+1}\) 到区间 \([0, b^{k+1} - 1]\) 的一个子集上的一个同构；但这个区间和 \(\mathbf{E}_{k+1}\) 具有相同数量的元素，即 \(b^{k+1}\) （第3号，命题5）；因此 \(f_{k+1}\) 是一个双射（§4，第2号，命题2，推论4），证明完成。

¶ 我们现在注意到，对于每个整数 \(a\) ，我们有 \(a < b^a\) 。这是通过对 \(a\) 进行归纳来证明的；对于 \(a = 0\) ，结果是显然的，而假设 \(a < b^a\) 蕴含 \(a + 1 \leqslant b^a < b \cdot b^a = b^{a+1}\) （第2号，命题3以及第4节，第2号，命题2）。因此存在一个最小整数 \(k\) 使得 \(a < b^k\) ，然后命题8表明存在唯一的有限序列 \((r_h)_{0 \leqslant h \leqslant k-1}\) ，使得 \(0 \leqslant r_h \leqslant b - 1\) （对于 \(0 \leqslant h \leqslant k - 1\) ），并且

\[
a = \sum_ {h = 0} ^ {k - 1} r _ {h} b ^ {k - h - 1}.
\]

此外，我们必须有 \(r_{0} > 0\) ，否则根据命题8， \(a < b^{k-1}\) 。表达式

\[
\sum_ {h = 0} ^ {k - 1} r _ {h} b ^ {k - h - 1}
\]

称为整数 a 的 b 进制展开。

* 在数学中所有不涉及数值计算的部分，命题 8 主要在应用于素数 \(b\) 时是有用的. *

当整数b足够小，以至于这样做切实可行时，我们可以用一个称为数字的独特符号来表示每个整数 \textless b 。表示0和1的数字通常是0和1。设a为一个整数。 \(_{k-1}\)

并令 \(\sum_{h=0}^{r_h b^{k-h-1}}\) 为其按基 \(b\) 的展开。如果整数 \(k\) （它

在此展开中出现）足够小以至于这样做可行，通常会将整数 \(a\) 与如下的符号序列相关联：把 \(r_0r_1 \ldots r_{k-2}r_{k-1}\) 从左到右书写，然后将每个整数 \(r_i\) 替换为代表它的数字；如此得到的符号称为与 \(a\) 相关联的数值符号。然后人们常常在包含它的项和关系中用其数值符号替换 \(a\) 。

例如，若C、Q、F、D是数字，则数值符号CQ、CQF、CQFD分别与 \(Cb + Q\) , \(Cb^{2} + Qb + F\) , \(Cb^{2} + Qb^{2} + Fb + D\) 相关联。

由命题8可知，与整数a相关联的数值符号是唯一的，并且如果 \(a < b^{k}\) ，它最多包含k位数字。请注意，与整数 \(b^{k}\) 相关联的数值符号由数字1后跟k个数字0组成。

这种用数字符号表示整数的方法称为以 b 为基数的记数系统。在实际数值计算中，使用以下系统：(a) 以 2 为基数的系统，即二进制系统，其中数字为 0 和 1；(b) 十进制系统，其中数字为 0、1、2、3、4、 \(5 = 4 + 1\) , \(6 = 5 + 1\) , \(7 = 6 + 1\) , \(8 = 7 + 1\) , \(9 = 8 + 1\) ，其中 \(b\) 是整数 \(9 + 1\) （其在此系统中的数字符号因此为 10）。

自中世纪以来，十进制系统在数值计算中一直被传统地使用，在本系列中，每当我们需要明确写出一个整数时，我们都将使用它。关于求两个由数字符号给出的整数的和、差、积及商的整数部分所对应的数字符号的方法，我们请读者参阅本系列中专门讨论数值计算的部分。

